\section{Cyclic obligations compiled to well-founded proofs}
\label{sec:cycles}

\subsection{A repeated goal is a candidate recursive call}

Existing tabling must not turn an active goal into its own proof. The proposed
extension preserves that rule. It changes what happens \emph{after} a recursive
pattern has been found: a cyclic dependency may be recorded as an unaccepted
proof-program candidate, whose recursive calls must subsequently be justified by
well-founded descent.

This direction is motivated by cyclic equational proof systems such as CycleQ,
and by recent work translating abstract cyclic proofs into induction
\cite{cycleq,unravelling}. The first Forge implementation should be substantially
more restricted than those general results. It should require a concrete scalar
or lexicographic ranking that Lean can use with ordinary well-founded induction.
There is no need to trust an arbitrary cyclic proof graph or implement a general
cyclic-proof kernel.

\subsection{The typed object to compile}

Each graph node $v$ has a telescope $\Gamma_v$, a state type $S_v$ packaging its
varying inputs, and a proposition family $P_v:S_v\to\mathsf{Prop}$. An edge
$e:v\to w$ carries a typed argument map $\sigma_e:S_v\to S_w$ together with
any branch condition under which that call is needed. A local constructor must
prove $P_v(x)$ from the recursive calls actually used in that branch and any
already checked nonrecursive facts.

These are logical edges, not just syntactic similarities. A matched variable
permutation needs a well-typed substitution; a cast between different index
types needs a proof; and a change to a dependent hypothesis needs an explicit
transport. The tagged state space is
\[
 S=\sum_{v:V}S_v,
 \qquad P(v,x)=P_v(x).
\]
The tag allows mutually recursive lemmas with different input types to be
compiled as one induction over a sum type. In an initial prototype integration,
a homogeneous telescope can simplify implementation, but it should not be baked
into the mathematical contract.

\begin{theorem}[Ranked local constructors]\label{thm:cyclic}
Let $\prec$ be a well-founded relation on $S$. Suppose that for every $s:S$,
a checked local constructor proves $P(s)$ from proofs of $P(t)$ for the recursive
calls it uses, and each such call has a checked proof $t\prec s$. Then
$\forall s,\,P(s)$.
\end{theorem}
\begin{proof}
Apply well-founded induction to $P$. At state $s$, the induction hypothesis
supplies $P(t)$ for each $t\prec s$. The checked descent proofs permit every
recursive call in the local constructor to be replaced by an instance of that
hypothesis. Applying the constructor proves $P(s)$.
\end{proof}

The theorem requires all branch cases to be covered and all local constructors
to be checked. A ranking of a graph alone does not meet those requirements.
The included Python code checks only the descent component for supplied graphs;
it does not extract those graphs from Lean or construct their logical branches.

\subsection{Why node phases matter}

Consider two obligations with calls
\[
 P(n)\longrightarrow Q(n),\qquad
 Q(n)\longrightarrow P(n-1)\quad(n\ge1).
\]
The first edge does not decrease $n$. Rejecting it independently would miss a
valid mutual-induction structure. A rank on the tagged state resolves this:
\[
 \rho(P,n)=2n+1,\qquad \rho(Q,n)=2n.
\]
Each edge now decreases by exactly one. The phase is not an arbitrary tie-breaker
in search order: it is part of the mathematical termination measure.

The prototype synthesizes this rank by enumerating small nonnegative weights
and node phases. A production solver can instead synthesize phases jointly with
affine weights. An acyclic run of zero-progress edges suggests phases cheaply:
assign reverse topological heights within a candidate phase graph, then scale
the primary measure sufficiently to dominate phase resets. Cycles of purely
zero-progress edges cannot be repaired by phases alone.

\subsection{Permuted arguments and lexicographic growth}

Two other fixtures distinguish useful ranking families.

For a call $(x,y)\mapsto(y,x-1)$ under $x\ge1$, neither a fixed-position
``first argument always decreases'' check nor a requirement that both arguments
shrink is appropriate. The rank $x+y$ decreases by one. This also covers the
standard subtractive Euclidean branches, when their positive-subtraction guards
are retained.

For calls
\[
 (x,y)\mapsto(x-1,2y)\quad(x\ge1),\qquad
 (x,y)\mapsto(x,y-1)\quad(y\ge1),
\]
lexicographic rank $(x,y)$ works. There is no globally decreasing nonnegative
affine scalar $a x+b y+c$ for both branches: the second requires $b>0$, while
the first would require $a-b y\ge1$ for unbounded $y$. This gives a meaningful
ablation between scalar and lexicographic ranking search.

The compiler must use lexicographic well-foundedness, not encode $(x,y)$ as
$a x+b y$ with a guessed fixed multiplier. Such an encoding is valid only when a
bound on the lower-priority component is available and proved.

\subsection{A tempting but invalid local rule}

Suppose the graph contains
\[
 (x,y)\mapsto(x-1,y+1)\quad(x\ge1),\qquad
 (x,y)\mapsto(x+1,y-1)\quad(y\ge1).
\]
Every edge decreases \emph{some} argument. Nevertheless,
\[
 (1,0)\longrightarrow(0,1)\longrightarrow(1,0)
\]
is an infinite repeatable cycle. The decrease on one edge is replenished by the
other. Neither a local decreasing-variable flag nor the existence of one
favorable cycle proves global well-foundedness.

The prototype's ranking search returns \code{unknown} here. Its size-change
closure also fails the acceptance criterion. The concrete two-cycle explains
this particular failure, but exhaustion of a ranking grammar in an arbitrary
case is not itself evidence of nontermination.

\subsection{Size-change closure as a proposal filter}

A size-change graph records relations between caller measures and callee measures.
The implementation uses entries $0,1,2$ for unknown, nonincrease, and strict
decrease. Along a composable path, an unknown comparison annihilates that path;
otherwise it is strict if either constituent is strict. Multiple paths between
an input and output retain the strongest established comparison.

Writing $\otimes$ for path composition and $\oplus$ for maximum strength,
\[
 (G\circ H)_{ij}=\bigoplus_k G_{ik}\otimes H_{kj}.
\]
Compute the closure of graphs of \emph{nonempty} paths, retaining source and
target node tags. The standard size-change test requires a strict diagonal
entry in every idempotent self-graph in that closure. A failure means that this
abstraction did not certify descent; a success relies on all recorded local
comparisons being semantically valid \cite{cycleq,unravelling}.

In this implementation the closure is diagnostic and guides ranking search.
It is not passed off as a Lean termination proof. The first compiler accepts
only an explicit checked ranking. A later release could formalize an abstract
size-change soundness theorem or implement a general unravelling construction,
but that is separate work with a separate proof obligation.

\subsection{Exact affine and lexicographic certificates}

The executable graph language uses natural-number measures and guarded affine
updates represented over the integers. For an edge $e:v\to w$ with update
$F_e$, a scalar ranking certificate proves
\[
 \rho_v(x)-\rho_w(F_e(x))-1\ge0.
\]
Subtracting one encodes strict decrease on integral ranks. Replacing this
condition by nonincrease is a soundness bug, not a harmless relaxation.

The checker proves the inequality by an explicit nonnegative rational
combination of the edge's guard polynomials, the natural-domain inequalities
$x_i\ge0$, and the constant inequality $1\ge0$. It separately proves that every
updated measure is nonnegative. Without the latter check, a syntactically
decreasing integer expression could leave the claimed well-founded domain.

For a lexicographic tuple, the certificate gives a pivot for every edge. All
components before the pivot must be equal, and the pivot decreases by at least
one. Later components need not decrease. The implemented checker requires
literal affine equality for the earlier components; a production extension may
accept guarded equality proved by two checked inequalities instead.

All ranking coefficients and phases in the prototype are nonnegative integers,
so the ranks map every natural state to a natural. More expressive ranks can
have signed coefficients, but then their nonnegativity over each node's invariant
region becomes an additional obligation. They should not be admitted merely
because they decrease along the sampled calls.

\subsection{A concrete production synthesis algorithm}

For one affine rank, use unknowns $w_{vi}$ and phases $b_v$ in
$\rho_v(x)=b_v+\sum_i w_{vi}x_i$. For each edge, introduce nonnegative Farkas
multipliers $\lambda_{ej}$ and equate coefficients in
\[
 \rho_v(x)-\rho_w(F_e(x))-1=\sum_j\lambda_{ej}g_{ej}(x).
\]
The update and guard coefficients are fixed, so this is a linear feasibility
problem in the rank parameters and multipliers. Solve over rationals, scale to
integral rank coefficients if necessary, regenerate all exact identities, and
check them. The solver's numerical feasibility verdict is not the certificate.

For a lexicographic tuple of fixed length, choose a pivot for each edge, enforce
prefix equality, and solve the resulting linear problem. Pivot choices can be
enumerated for small SCCs or represented by finite selectors in an SMT/MILP
proposer. The backend must export the selected pivots and exact coefficient
witnesses, not a floating-point tuple that is rounded without rechecking.

The Python implementation uses bounded enumeration instead: common weights
across nodes, small individual phases, then permutations of the original
coordinates as a lexicographic fallback. Cone membership is solved by enumerating
linearly independent supports of size at most the affine coefficient dimension.
This is appropriate for its one- and two-variable tests; it is not proposed as
a scalable replacement for an exact LP backend.

\subsection{Compiling an SCC without introducing circular declarations}

A practical Lean compiler can close one strongly connected component as follows.
First, package its node goals into the tagged family $P$. Second, construct a
well-founded relation from the accepted rank. Third, build one local proof
constructor per node using ordinary case splits and the already checked
nonrecursive facts. Recursive edges become accesses to the well-founded
induction hypothesis, instantiated at the typed callee state and supplied with
the corresponding descent proof. Finally, project the theorem for the entry
node and original arguments.

No node theorem should be inserted into the environment as an assumed theorem
while its SCC is being checked. The only recursive authority is the induction
hypothesis provided by the well-founded recursor. Lean supports structural and
well-founded recursive definitions through ordinary kernel-checked terms;
\code{termination\_by} and \code{decreasing\_by} are possible elaboration
interfaces, not a new trusted proof rule \cite{lean-recursion,lean-inductive}.

The first implementation should support explicit tagged states and \emph{given}
local constructors, then add graph extraction. Trying to discover dependent
back-edge substitutions, generate branches, infer ranks, and lower proofs all
at once makes failures difficult to localize. The supplied examples provide
separate descent gates before that larger integration begins.
